% Topic content; common layout lives in the shared course preamble.
\newcommand{\CourseTemplatePath}{../../../../../templates/slides/}
\input{\CourseTemplatePath course_preamble.tex}
\title[Q Functions, Optimality, and VFI]{Lecture: Q Functions, Optimality,\newline and Value Function Iteration}
\newcommand{\E}{\mathbb E}
\newcommand{\Sset}{\mathcal S}
\newcommand{\Aset}{\mathcal A}
\tikzset{valueState/.style={draw=SessionTeal,fill=SessionTealLight,rounded corners=1mm,minimum width=15mm,minimum height=9mm,align=center},valueFlow/.style={-{Stealth[length=2mm]},draw=SessionCharcoal,line width=.8pt},stateCircle/.style={circle,draw=SessionTeal,fill=white,minimum size=6mm,inner sep=1pt},actionCircle/.style={circle,draw=SessionTeal,fill=SessionTeal,minimum size=3mm,inner sep=0pt}}
% The maze geometry and reward convention match the Week 1 example.
\definecolor{ValueMazeWall}{HTML}{928580}
\newcommand{\ValueMazeBase}{%
  \fill[white] (0,0) rectangle (9,6);
  \foreach \x/\y in {2/2,2/3,2/4,7/3,7/4,7/5,5/1}
    \fill[ValueMazeWall] (\x,\y) rectangle +(1,1);
  \draw[SessionLine,line width=.3pt,step=1] (0,0) grid (9,6);
  \draw[SessionCharcoal,line width=.7pt] (0,0) rectangle (9,6);
  % S and G label the endpoint states; other free cells are numbered.
  \node[anchor=north west,inner sep=0pt,font=\sffamily\fontsize{6}{6}\selectfont,text=SessionMuted!65] at (0.09,5.92) {1};
  \node[anchor=north west,inner sep=0pt,font=\sffamily\fontsize{6}{6}\selectfont,text=SessionMuted!65] at (1.09,5.92) {2};
  \node[anchor=north west,inner sep=0pt,font=\sffamily\fontsize{6}{6}\selectfont,text=SessionMuted!65] at (2.09,5.92) {3};
  \node[anchor=north west,inner sep=0pt,font=\sffamily\fontsize{6}{6}\selectfont,text=SessionMuted!65] at (3.09,5.92) {4};
  \node[anchor=north west,inner sep=0pt,font=\sffamily\fontsize{6}{6}\selectfont,text=SessionMuted!65] at (4.09,5.92) {5};
  \node[anchor=north west,inner sep=0pt,font=\sffamily\fontsize{6}{6}\selectfont,text=SessionMuted!65] at (5.09,5.92) {6};
  \node[anchor=north west,inner sep=0pt,font=\sffamily\fontsize{6}{6}\selectfont,text=SessionMuted!65] at (6.09,5.92) {7};
  \node[anchor=north west,inner sep=0pt,font=\sffamily\fontsize{6}{6}\selectfont,text=SessionMuted!65] at (0.09,4.92) {8};
  \node[anchor=north west,inner sep=0pt,font=\sffamily\fontsize{6}{6}\selectfont,text=SessionMuted!65] at (1.09,4.92) {9};
  \node[anchor=north west,inner sep=0pt,font=\sffamily\fontsize{6}{6}\selectfont,text=SessionMuted!65] at (3.09,4.92) {10};
  \node[anchor=north west,inner sep=0pt,font=\sffamily\fontsize{6}{6}\selectfont,text=SessionMuted!65] at (4.09,4.92) {11};
  \node[anchor=north west,inner sep=0pt,font=\sffamily\fontsize{6}{6}\selectfont,text=SessionMuted!65] at (5.09,4.92) {12};
  \node[anchor=north west,inner sep=0pt,font=\sffamily\fontsize{6}{6}\selectfont,text=SessionMuted!65] at (6.09,4.92) {13};
  \node[anchor=north west,inner sep=0pt,font=\sffamily\fontsize{6}{6}\selectfont,text=SessionMuted!65] at (8.09,4.92) {14};
  \node[anchor=north west,inner sep=0pt,font=\sffamily\fontsize{6}{6}\selectfont,text=SessionMuted!65] at (1.09,3.92) {15};
  \node[anchor=north west,inner sep=0pt,font=\sffamily\fontsize{6}{6}\selectfont,text=SessionMuted!65] at (3.09,3.92) {16};
  \node[anchor=north west,inner sep=0pt,font=\sffamily\fontsize{6}{6}\selectfont,text=SessionMuted!65] at (4.09,3.92) {17};
  \node[anchor=north west,inner sep=0pt,font=\sffamily\fontsize{6}{6}\selectfont,text=SessionMuted!65] at (5.09,3.92) {18};
  \node[anchor=north west,inner sep=0pt,font=\sffamily\fontsize{6}{6}\selectfont,text=SessionMuted!65] at (6.09,3.92) {19};
  \node[anchor=north west,inner sep=0pt,font=\sffamily\fontsize{6}{6}\selectfont,text=SessionMuted!65] at (8.09,3.92) {20};
  \node[anchor=north west,inner sep=0pt,font=\sffamily\fontsize{6}{6}\selectfont,text=SessionMuted!65] at (0.09,2.92) {21};
  \node[anchor=north west,inner sep=0pt,font=\sffamily\fontsize{6}{6}\selectfont,text=SessionMuted!65] at (1.09,2.92) {22};
  \node[anchor=north west,inner sep=0pt,font=\sffamily\fontsize{6}{6}\selectfont,text=SessionMuted!65] at (3.09,2.92) {23};
  \node[anchor=north west,inner sep=0pt,font=\sffamily\fontsize{6}{6}\selectfont,text=SessionMuted!65] at (4.09,2.92) {24};
  \node[anchor=north west,inner sep=0pt,font=\sffamily\fontsize{6}{6}\selectfont,text=SessionMuted!65] at (5.09,2.92) {25};
  \node[anchor=north west,inner sep=0pt,font=\sffamily\fontsize{6}{6}\selectfont,text=SessionMuted!65] at (6.09,2.92) {26};
  \node[anchor=north west,inner sep=0pt,font=\sffamily\fontsize{6}{6}\selectfont,text=SessionMuted!65] at (7.09,2.92) {27};
  \node[anchor=north west,inner sep=0pt,font=\sffamily\fontsize{6}{6}\selectfont,text=SessionMuted!65] at (8.09,2.92) {28};
  \node[anchor=north west,inner sep=0pt,font=\sffamily\fontsize{6}{6}\selectfont,text=SessionMuted!65] at (0.09,1.92) {29};
  \node[anchor=north west,inner sep=0pt,font=\sffamily\fontsize{6}{6}\selectfont,text=SessionMuted!65] at (1.09,1.92) {30};
  \node[anchor=north west,inner sep=0pt,font=\sffamily\fontsize{6}{6}\selectfont,text=SessionMuted!65] at (2.09,1.92) {31};
  \node[anchor=north west,inner sep=0pt,font=\sffamily\fontsize{6}{6}\selectfont,text=SessionMuted!65] at (3.09,1.92) {32};
  \node[anchor=north west,inner sep=0pt,font=\sffamily\fontsize{6}{6}\selectfont,text=SessionMuted!65] at (4.09,1.92) {33};
  \node[anchor=north west,inner sep=0pt,font=\sffamily\fontsize{6}{6}\selectfont,text=SessionMuted!65] at (6.09,1.92) {34};
  \node[anchor=north west,inner sep=0pt,font=\sffamily\fontsize{6}{6}\selectfont,text=SessionMuted!65] at (7.09,1.92) {35};
  \node[anchor=north west,inner sep=0pt,font=\sffamily\fontsize{6}{6}\selectfont,text=SessionMuted!65] at (8.09,1.92) {36};
  \node[anchor=north west,inner sep=0pt,font=\sffamily\fontsize{6}{6}\selectfont,text=SessionMuted!65] at (0.09,0.92) {37};
  \node[anchor=north west,inner sep=0pt,font=\sffamily\fontsize{6}{6}\selectfont,text=SessionMuted!65] at (1.09,0.92) {38};
  \node[anchor=north west,inner sep=0pt,font=\sffamily\fontsize{6}{6}\selectfont,text=SessionMuted!65] at (2.09,0.92) {39};
  \node[anchor=north west,inner sep=0pt,font=\sffamily\fontsize{6}{6}\selectfont,text=SessionMuted!65] at (3.09,0.92) {40};
  \node[anchor=north west,inner sep=0pt,font=\sffamily\fontsize{6}{6}\selectfont,text=SessionMuted!65] at (4.09,0.92) {41};
  \node[anchor=north west,inner sep=0pt,font=\sffamily\fontsize{6}{6}\selectfont,text=SessionMuted!65] at (5.09,0.92) {42};
  \node[anchor=north west,inner sep=0pt,font=\sffamily\fontsize{6}{6}\selectfont,text=SessionMuted!65] at (6.09,0.92) {43};
  \node[anchor=north west,inner sep=0pt,font=\sffamily\fontsize{6}{6}\selectfont,text=SessionMuted!65] at (7.09,0.92) {44};
  \node[anchor=north west,inner sep=0pt,font=\sffamily\fontsize{6}{6}\selectfont,text=SessionMuted!65] at (8.09,0.92) {45};
}
\newcommand{\ValueMazeLabels}{%
  \node[font=\sffamily\large,text=SessionCharcoal,fill=white,inner sep=0pt] at (.5,3.5) {S};
  \node[font=\sffamily\large,text=SessionCharcoal,fill=white,inner sep=0pt] at (8.5,5.5) {G};
}


% Concrete deterministic policy and its undiscounted, terminating values.
\newcommand{\MazePolicy}{%
\node[anchor=north west,inner sep=0pt,font=\sffamily\fontsize{6}{6}\selectfont,text=SessionMuted!65] at (0.09,3.92) {S};
\node[anchor=north west,inner sep=0pt,font=\sffamily\fontsize{6}{6}\selectfont,text=SessionMuted!65] at (8.09,5.92) {G};
\node[font=\sffamily\large,text=SessionCharcoal,inner sep=0pt] at (8.5,5.5) {U};
\node[font=\sffamily\large,text=SessionCharcoal,inner sep=0pt] at (0.5,5.5) {R};
\node[font=\sffamily\large,text=SessionCharcoal,inner sep=0pt] at (1.5,5.5) {R};
\node[font=\sffamily\large,text=SessionCharcoal,inner sep=0pt] at (2.5,5.5) {R};
\node[font=\sffamily\large,text=SessionCharcoal,inner sep=0pt] at (3.5,5.5) {D};
\node[font=\sffamily\large,text=SessionCharcoal,inner sep=0pt] at (4.5,5.5) {D};
\node[font=\sffamily\large,text=SessionCharcoal,inner sep=0pt] at (5.5,5.5) {D};
\node[font=\sffamily\large,text=SessionCharcoal,inner sep=0pt] at (6.5,5.5) {D};
\node[font=\sffamily\large,text=SessionCharcoal,inner sep=0pt] at (0.5,4.5) {U};
\node[font=\sffamily\large,text=SessionCharcoal,inner sep=0pt] at (1.5,4.5) {U};
\node[font=\sffamily\large,text=SessionCharcoal,inner sep=0pt] at (3.5,4.5) {D};
\node[font=\sffamily\large,text=SessionCharcoal,inner sep=0pt] at (4.5,4.5) {D};
\node[font=\sffamily\large,text=SessionCharcoal,inner sep=0pt] at (5.5,4.5) {D};
\node[font=\sffamily\large,text=SessionCharcoal,inner sep=0pt] at (6.5,4.5) {D};
\node[font=\sffamily\large,text=SessionCharcoal,inner sep=0pt] at (8.5,4.5) {U};
\node[font=\sffamily\large,text=SessionCharcoal,inner sep=0pt] at (0.5,3.5) {D};
\node[font=\sffamily\large,text=SessionCharcoal,inner sep=0pt] at (1.5,3.5) {D};
\node[font=\sffamily\large,text=SessionCharcoal,inner sep=0pt] at (3.5,3.5) {R};
\node[font=\sffamily\large,text=SessionCharcoal,inner sep=0pt] at (4.5,3.5) {R};
\node[font=\sffamily\large,text=SessionCharcoal,inner sep=0pt] at (5.5,3.5) {R};
\node[font=\sffamily\large,text=SessionCharcoal,inner sep=0pt] at (6.5,3.5) {D};
\node[font=\sffamily\large,text=SessionCharcoal,inner sep=0pt] at (8.5,3.5) {U};
\node[font=\sffamily\large,text=SessionCharcoal,inner sep=0pt] at (0.5,2.5) {D};
\node[font=\sffamily\large,text=SessionCharcoal,inner sep=0pt] at (1.5,2.5) {D};
\node[font=\sffamily\large,text=SessionCharcoal,inner sep=0pt] at (3.5,2.5) {U};
\node[font=\sffamily\large,text=SessionCharcoal,inner sep=0pt] at (4.5,2.5) {R};
\node[font=\sffamily\large,text=SessionCharcoal,inner sep=0pt] at (5.5,2.5) {R};
\node[font=\sffamily\large,text=SessionCharcoal,inner sep=0pt] at (6.5,2.5) {R};
\node[font=\sffamily\large,text=SessionCharcoal,inner sep=0pt] at (7.5,2.5) {R};
\node[font=\sffamily\large,text=SessionCharcoal,inner sep=0pt] at (8.5,2.5) {U};
\node[font=\sffamily\large,text=SessionCharcoal,inner sep=0pt] at (0.5,1.5) {R};
\node[font=\sffamily\large,text=SessionCharcoal,inner sep=0pt] at (1.5,1.5) {R};
\node[font=\sffamily\large,text=SessionCharcoal,inner sep=0pt] at (2.5,1.5) {R};
\node[font=\sffamily\large,text=SessionCharcoal,inner sep=0pt] at (3.5,1.5) {U};
\node[font=\sffamily\large,text=SessionCharcoal,inner sep=0pt] at (4.5,1.5) {U};
\node[font=\sffamily\large,text=SessionCharcoal,inner sep=0pt] at (6.5,1.5) {U};
\node[font=\sffamily\large,text=SessionCharcoal,inner sep=0pt] at (7.5,1.5) {U};
\node[font=\sffamily\large,text=SessionCharcoal,inner sep=0pt] at (8.5,1.5) {U};
\node[font=\sffamily\large,text=SessionCharcoal,inner sep=0pt] at (0.5,0.5) {U};
\node[font=\sffamily\large,text=SessionCharcoal,inner sep=0pt] at (1.5,0.5) {U};
\node[font=\sffamily\large,text=SessionCharcoal,inner sep=0pt] at (2.5,0.5) {U};
\node[font=\sffamily\large,text=SessionCharcoal,inner sep=0pt] at (3.5,0.5) {U};
\node[font=\sffamily\large,text=SessionCharcoal,inner sep=0pt] at (4.5,0.5) {U};
\node[font=\sffamily\large,text=SessionCharcoal,inner sep=0pt] at (5.5,0.5) {R};
\node[font=\sffamily\large,text=SessionCharcoal,inner sep=0pt] at (6.5,0.5) {U};
\node[font=\sffamily\large,text=SessionCharcoal,inner sep=0pt] at (7.5,0.5) {U};
\node[font=\sffamily\large,text=SessionCharcoal,inner sep=0pt] at (8.5,0.5) {U};
}
\newcommand{\MazePolicyValues}{%
\node[anchor=north west,inner sep=0pt,font=\sffamily\fontsize{6}{6}\selectfont,text=SessionMuted!65] at (0.09,3.92) {S};
\node[anchor=north west,inner sep=0pt,font=\sffamily\fontsize{6}{6}\selectfont,text=SessionMuted!65] at (8.09,5.92) {G};
\node[font=\sffamily\large,text=SessionCharcoal,inner sep=0pt] at (8.5,5.5) {0};
\node[font=\sffamily\large,text=SessionCharcoal,inner sep=0pt] at (0.5,5.5) {-14};
\node[font=\sffamily\large,text=SessionCharcoal,inner sep=0pt] at (1.5,5.5) {-13};
\node[font=\sffamily\large,text=SessionCharcoal,inner sep=0pt] at (2.5,5.5) {-12};
\node[font=\sffamily\large,text=SessionCharcoal,inner sep=0pt] at (3.5,5.5) {-11};
\node[font=\sffamily\large,text=SessionCharcoal,inner sep=0pt] at (4.5,5.5) {-10};
\node[font=\sffamily\large,text=SessionCharcoal,inner sep=0pt] at (5.5,5.5) {-9};
\node[font=\sffamily\large,text=SessionCharcoal,inner sep=0pt] at (6.5,5.5) {-8};
\node[font=\sffamily\large,text=SessionCharcoal,inner sep=0pt] at (0.5,4.5) {-15};
\node[font=\sffamily\large,text=SessionCharcoal,inner sep=0pt] at (1.5,4.5) {-14};
\node[font=\sffamily\large,text=SessionCharcoal,inner sep=0pt] at (3.5,4.5) {-10};
\node[font=\sffamily\large,text=SessionCharcoal,inner sep=0pt] at (4.5,4.5) {-9};
\node[font=\sffamily\large,text=SessionCharcoal,inner sep=0pt] at (5.5,4.5) {-8};
\node[font=\sffamily\large,text=SessionCharcoal,inner sep=0pt] at (6.5,4.5) {-7};
\node[font=\sffamily\large,text=SessionCharcoal,inner sep=0pt] at (8.5,4.5) {-1};
\node[font=\sffamily\large,text=SessionCharcoal,inner sep=0pt] at (0.5,3.5) {-16};
\node[font=\sffamily\large,text=SessionCharcoal,inner sep=0pt] at (1.5,3.5) {-15};
\node[font=\sffamily\large,text=SessionCharcoal,inner sep=0pt] at (3.5,3.5) {-9};
\node[font=\sffamily\large,text=SessionCharcoal,inner sep=0pt] at (4.5,3.5) {-8};
\node[font=\sffamily\large,text=SessionCharcoal,inner sep=0pt] at (5.5,3.5) {-7};
\node[font=\sffamily\large,text=SessionCharcoal,inner sep=0pt] at (6.5,3.5) {-6};
\node[font=\sffamily\large,text=SessionCharcoal,inner sep=0pt] at (8.5,3.5) {-2};
\node[font=\sffamily\large,text=SessionCharcoal,inner sep=0pt] at (0.5,2.5) {-15};
\node[font=\sffamily\large,text=SessionCharcoal,inner sep=0pt] at (1.5,2.5) {-14};
\node[font=\sffamily\large,text=SessionCharcoal,inner sep=0pt] at (3.5,2.5) {-10};
\node[font=\sffamily\large,text=SessionCharcoal,inner sep=0pt] at (4.5,2.5) {-7};
\node[font=\sffamily\large,text=SessionCharcoal,inner sep=0pt] at (5.5,2.5) {-6};
\node[font=\sffamily\large,text=SessionCharcoal,inner sep=0pt] at (6.5,2.5) {-5};
\node[font=\sffamily\large,text=SessionCharcoal,inner sep=0pt] at (7.5,2.5) {-4};
\node[font=\sffamily\large,text=SessionCharcoal,inner sep=0pt] at (8.5,2.5) {-3};
\node[font=\sffamily\large,text=SessionCharcoal,inner sep=0pt] at (0.5,1.5) {-14};
\node[font=\sffamily\large,text=SessionCharcoal,inner sep=0pt] at (1.5,1.5) {-13};
\node[font=\sffamily\large,text=SessionCharcoal,inner sep=0pt] at (2.5,1.5) {-12};
\node[font=\sffamily\large,text=SessionCharcoal,inner sep=0pt] at (3.5,1.5) {-11};
\node[font=\sffamily\large,text=SessionCharcoal,inner sep=0pt] at (4.5,1.5) {-8};
\node[font=\sffamily\large,text=SessionCharcoal,inner sep=0pt] at (6.5,1.5) {-6};
\node[font=\sffamily\large,text=SessionCharcoal,inner sep=0pt] at (7.5,1.5) {-5};
\node[font=\sffamily\large,text=SessionCharcoal,inner sep=0pt] at (8.5,1.5) {-4};
\node[font=\sffamily\large,text=SessionCharcoal,inner sep=0pt] at (0.5,0.5) {-15};
\node[font=\sffamily\large,text=SessionCharcoal,inner sep=0pt] at (1.5,0.5) {-14};
\node[font=\sffamily\large,text=SessionCharcoal,inner sep=0pt] at (2.5,0.5) {-13};
\node[font=\sffamily\large,text=SessionCharcoal,inner sep=0pt] at (3.5,0.5) {-12};
\node[font=\sffamily\large,text=SessionCharcoal,inner sep=0pt] at (4.5,0.5) {-9};
\node[font=\sffamily\large,text=SessionCharcoal,inner sep=0pt] at (5.5,0.5) {-8};
\node[font=\sffamily\large,text=SessionCharcoal,inner sep=0pt] at (6.5,0.5) {-7};
\node[font=\sffamily\large,text=SessionCharcoal,inner sep=0pt] at (7.5,0.5) {-6};
\node[font=\sffamily\large,text=SessionCharcoal,inner sep=0pt] at (8.5,0.5) {-5};
}

\begin{document}
\makecoursetitle

\begin{frame}{Lecture Overview}
\begin{itemize}
\item \textbf{Action Values: The Q Function}
  \begin{itemize}\item Fixing the first action\item How \(Q^\sigma\) and \(V^\sigma\) are related\end{itemize}
\item \textbf{The Principle of Optimality}
  \begin{itemize}\item Optimal continuation from the next state\item Bellman optimality equations and branching diagrams\end{itemize}
\item \textbf{Value Function Iteration}
  \begin{itemize}\item The Bellman update, convergence, and a stopping rule\item Recovering a policy; job search and continuous saving\end{itemize}
\end{itemize}
\end{frame}

\begin{frame}{Example: A Deterministic Maze Policy}
\small
A fixed policy \(\sigma\) selects one action in every free cell. Set \(\beta=1\) for this terminating example.
\begin{columns}[c]
\begin{column}{.60\textwidth}
\begin{figure}\centering
\begin{tikzpicture}[x=.85cm,y=.85cm]
\ValueMazeBase\MazePolicy
\end{tikzpicture}
\caption{The action selected by \(\sigma\).}
\end{figure}
\end{column}
\begin{column}{.37\textwidth}
\begin{itemize}
\item U: up; D: down;\\L: left; R: right.
\item Each move gives \(-1\), including the move into G.
\item G is absorbing, with reward zero. Its displayed U keeps the process at G.
\item Following this policy reaches G from every free cell.
\end{itemize}
\end{column}\end{columns}
\end{frame}

\begin{frame}{Example: Computing the Policy's Values}
\small
With \(\beta=1\), compute values backward from G, following the actions of \(\sigma\).
\begin{columns}[c]
\begin{column}{.60\textwidth}
\begin{figure}\centering
\begin{tikzpicture}[x=.85cm,y=.85cm]
\ValueMazeBase\MazePolicyValues
\end{tikzpicture}
\caption{State values under the same policy.}
\end{figure}
\end{column}
\begin{column}{.37\textwidth}
\(V^\sigma(\mathrm G)=0\). For every other cell,
\[V^\sigma(s)=-1+V^\sigma(s'),\]
where \(s'\) is its chosen next cell.
\begin{align*}
V^\sigma(14)&=-1+0=-1,\\
V^\sigma(20)&=-1+(-1)=-2,\\
V^\sigma(28)&=-1+(-2)=-3.
\end{align*}
Continuing backward gives \(V^\sigma(\mathrm S)=-16\).
\par\smallskip
Each value is minus the number of moves to G under \(\sigma\).
\end{column}\end{columns}
\end{frame}

\begin{frame}{The Action-Value Function}
The \term{action-value function} \(Q^\sigma(s,a)\) is the expected return when we start at \(s\), take action \(a\), and follow \(\sigma\) thereafter:
\[Q^\sigma(s,a)=\E_\sigma\!\left[\sum_{\tau=1}^{\infty}\beta^{\tau-1}r_{t+\tau}
\,\middle|\,s_t=s,\ a_t=a\right].\]
\begin{itemize}
\item The first action \(a\) is \textbf{fixed}; actions from date \(t+1\) onward follow \(\sigma\).
\item \(a\) may differ from the action that \(\sigma\) would choose at \(s\). Interpret the conditioning as deliberately taking \(a\), even if \(\sigma\) never selects it.
\item Rewards and subsequent states can still be random.
\end{itemize}
\end{frame}

\begin{frame}{Linking Q and V}
\small
Fixing the first action gives
\[Q^\sigma(s,a)=\E\!\left[r_{t+1}+\beta V^\sigma(s_{t+1})\mid s_t=s,a_t=a\right].\]
For discrete rewards and next states, expand this expectation:
\[Q^\sigma(s,a)=\sum_r\Pr\{R(s,a)=r\}\sum_{s'}P(s'\mid s,a,r)
\bigl[r+\beta V^\sigma(s')\bigr].\]
Averaging the first action according to the policy recovers the state value:
\[V^\sigma(s)=\sum_{a\in\Aset(s)}\Pr\{\sigma(s)=a\}\,Q^\sigma(s,a).\]
For a deterministic policy, \(V^\sigma(s)=Q^\sigma(s,\sigma(s))\).
\par\smallskip
\footnotesize The conditional kernel allows reward--next-state dependence. Replace sums by probability integrals for continuous outcomes.
\end{frame}

\begin{frame}{Why Use Q?}
\begin{itemize}
\item \textbf{Compare actions at the same state.} \(Q^\sigma(s,a)\) includes the immediate reward and the consequences of following \(\sigma\) afterward.
\item \textbf{Choose an action directly.} Given action values, a \term{greedy action} is
\[a_{\mathrm g}(s)\in\argmax_{a\in\Aset(s)}Q^\sigma(s,a).\]
\item \textbf{Learn from experience.} Returns from rollouts starting with \((s,a)\) can estimate \(Q^\sigma(s,a)\); an explicit transition model is not required for this estimation.
\end{itemize}
\small A greedy choice from \(Q^\sigma\) uses \(\sigma\)'s continuation values. It need not already be optimal. Storing action values also takes more space than storing one value per state.
\end{frame}

\begin{frame}{Example: An Optimal Continuation}
\small
Keep the route from S to cell 23 fixed. At cell 23, compare two continuations to G.
\begin{columns}[c]
\begin{column}{.56\textwidth}
\begin{figure}\centering
\begin{tikzpicture}[x=.62cm,y=.62cm]
\ValueMazeBase
\foreach \x/\y/\u/\v in {0/3/0/2,0/2/0/1,0/1/1/1,1/1/2/1,2/1/3/1,3/1/3/2}
 \draw[valueFlow,shorten <=2pt,shorten >=2pt] (\x+.5,\y+.5)--(\u+.5,\v+.5);
\foreach \x/\y/\u/\v in {3/2/4/2,4/2/5/2,5/2/6/2,6/2/7/2,7/2/8/2,8/2/8/3,8/3/8/4,8/4/8/5}
 \draw[valueFlow,draw=SessionTeal,shorten <=2pt,shorten >=2pt] (\x+.5,\y+.5)--(\u+.5,\v+.5);
\foreach \x/\y/\u/\v in {3/2/3/3,3/3/4/3,4/3/5/3,5/3/6/3,6/3/6/2}
 \draw[valueFlow,draw=SessionRed,dashed,shorten <=2pt,shorten >=2pt] (\x+.5,\y+.5)--(\u+.5,\v+.5);
\draw[CourseVariableColor,line width=1pt] (3.5,2.5) circle (.3);
\ValueMazeLabels
\end{tikzpicture}
\caption{A common prefix and two continuations.}
\end{figure}
\end{column}
\begin{column}{.41\textwidth}
\begin{itemize}
\item \textcolor{SessionTeal}{Solid continuation:} 8 moves.
\item \textcolor{SessionRed}{Dashed detour:} 10 moves, including the shared final segment.
\end{itemize}
At cell 23, a continuation of \(L\) moves gives
\[-L\qquad(\beta=1).\]
Replacing the detour raises \(V(\mathrm S)\) from \(-16\) to \(-14\): a gain of 2.
\end{column}\end{columns}
\end{frame}

\begin{frame}{The Principle of Optimality}
For an optimal policy, the continuation must be optimal from each state reached with positive probability.
\begin{itemize}
\item Suppose a better continuation is available at a reached state.
\item Keep all earlier choices unchanged and switch to that better continuation.
\item The original plan's expected return then increases, contradicting its optimality.
\end{itemize}
\par\medskip
The \term{Markov state} contains the information needed for future decisions. Earlier rewards are already realized; they do not change which continuation is best.
\par\smallskip
\small With uncertainty, the continuation is a policy for every possible next state. Optimality concerns \textbf{expected return}, not whether each realized path is the shortest or luckiest.
\end{frame}

\begin{frame}{Optimal State and Action Values}
For a finite MDP with bounded rewards and \(0<\beta<1\), an optimal stationary deterministic policy exists. Define
\[V^*(s)=\max_\sigma V^\sigma(s),\qquad
Q^*(s,a)=\max_\sigma Q^\sigma(s,a).\]
The principle of optimality gives
\begin{align*}
Q^*(s,a)&=\E\!\left[r_{t+1}+\beta V^*(s_{t+1})\mid s_t=s,a_t=a\right],\\
V^*(s)&=\max_{a\in\Aset(s)}Q^*(s,a).
\end{align*}
\begin{itemize}
\item \(Q^*\): fix the current action, then choose optimally afterward.
\item \(V^*\): also choose the current action optimally.
\item The expectation averages the outcomes of the chosen action; it does not let us choose the random outcome.
\end{itemize}
\end{frame}

\begin{frame}{Optimality as a One-Step Look Ahead}
\small
\begin{columns}[T]
\begin{column}{.49\textwidth}\centering
\(V^*(s)=\max_a Q^*(s,a)\)
\begin{figure}\centering
\begin{tikzpicture}[x=1cm,y=1cm,font=\sffamily\footnotesize]
\node[stateCircle] (s) at (0,0) {\(s\)};
\foreach \i/\x in {1/-1.35,2/1.35}{
 \node[actionCircle,label={right:\(a_\i\)}] (a\i) at (\x,-1) {};
 \draw[valueFlow] (s)--(a\i);
 \foreach \j/\dx in {1/-.55,2/.55}{
  \node[stateCircle] (t\i\j) at (\x+\dx,-2.2) {};
  \draw[valueFlow] (a\i)--(t\i\j);
 }
}
\node[text=SessionTeal] at (0,-.55) {maximize};
\node[text=SessionMuted] at (0,-2.85) {average \(r+\beta V^*(s')\) for each action};
\end{tikzpicture}
\caption{Choose now, average outcomes.}
\end{figure}\end{column}
\begin{column}{.49\textwidth}\centering
\(Q^*(s,a)=\E[r+\beta\max_{a'}Q^*(s',a')\mid s,a]\)
\begin{figure}\centering
\begin{tikzpicture}[x=1cm,y=1cm,font=\sffamily\footnotesize]
\node[actionCircle,label={above:\((s,a)\)}] (a) at (0,0) {};
\foreach \i/\x in {1/-1.35,2/1.35}{
 \node[stateCircle] (s\i) at (\x,-1) {\(s'_\i\)};
 \draw[valueFlow] (a)--(s\i);
 \foreach \j/\dx in {1/-.55,2/.55}{
  \node[actionCircle] (b\i\j) at (\x+\dx,-2.2) {};
  \draw[valueFlow] (s\i)--(b\i\j);
 }
 \node[text=SessionTeal] at (\x,-1.72) {maximize};
}
\node[text=SessionMuted] at (0,-.4) {average};
\node[text=SessionMuted] at (0,-2.85) {choose \(a'\) after observing \(s'\)};
\end{tikzpicture}
\caption{Fix now, choose at the next state.}
\end{figure}\end{column}\end{columns}
\footnotesize Open circles: states. Filled circles: state--action pairs. Adapted from Sutton and Barto (2018), Figure 3.4; branches group reward and next-state outcomes.
\end{frame}

\begin{frame}{The Bellman Optimality Operator}
An \term{operator} maps a candidate value function \(v\) to a new function:
\[(Tv)(s)=\max_{a\in\Aset(s)}
\sum_r\Pr\{R(s,a)=r\}\sum_{s'}P(s'\mid s,a,r)
\bigl[r+\beta v(s')\bigr].\]
\begin{itemize}
\item \(v(s')\) is our current estimate of the continuation value.
\item For each action, average current reward plus discounted continuation value; retain the largest result.
\item The optimal value is a \term{fixed point}: \(TV^*=V^*\).
\item Policy evaluation keeps a given policy fixed. This operator maximizes over actions at every update.
\end{itemize}
\small The sums describe discrete outcomes; probability integrals give the continuous version.
\end{frame}

\begin{frame}{Value Function Iteration}
\small
For a finite MDP with known reward and transition distributions:
\begin{enumerate}
\item Initialize \(v_0(s)=0\) for all states; keep absorbing zero-reward terminal states at zero.
\item For every state \(s\), score each feasible action using the \textbf{old} function:
\[q_k(s,a)=\sum_r\Pr\{R(s,a)=r\}\sum_{s'}P(s'\mid s,a,r)
\bigl[r+\beta v_k(s')\bigr].\]
\item Store the best score in a new function: \(v_{k+1}(s)=\max_a q_k(s,a)\).
\item Check the stopping rule; otherwise replace \(v_k\) by \(v_{k+1}\) and repeat.
\end{enumerate}
Here \(q_k\) is an action score based on the current approximation; it need not equal \(Q^\sigma\) for a particular policy.
\par\smallskip
One \term{sweep} updates every state once. Using a separate new array makes these updates synchronous.
\end{frame}

\begin{frame}{Example: Values Move Back Through the Maze}
\small
Consider cells 28, 20, and 14 near G in the same maze. Set \(\beta=0.9\), \(v_0=0\); each move before G costs one.
\begin{figure}\centering
\begin{tikzpicture}[font=\sffamily\small]
\foreach \i/\lab in {0/28,1/20,2/14,3/G}
 \node[valueState] (s\i) at (2.6*\i,0) {\lab};
\foreach \i/\j in {0/1,1/2,2/3}
 \draw[valueFlow] (s\i)--node[above]{up}(s\j);
\end{tikzpicture}
\caption{The final corridor to G.}
\end{figure}
\begin{center}
\begin{tabular}{lrrrr}\toprule
 & Cell 28 & Cell 20 & Cell 14 & G\\\midrule
\(v_0\)&0&0&0&0\\
\(v_1\)&\(-1\)&\(-1\)&\(-1\)&0\\
\(v_2\)&\(-1.9\)&\(-1.9\)&\(-1\)&0\\
\(v_3\)&\(-2.71\)&\(-1.9\)&\(-1\)&0\\\bottomrule
\end{tabular}
\end{center}
For example, \(v_3(28)=-1+0.9v_2(20)=-2.71\).
With zero initialization, \(v_k\) is the optimal \(k\)-period return with zero value afterward.
\end{frame}

\begin{frame}{Why VFI Converges}
\small
For finite states, use \(\|v-w\|_\infty=\max_s|v(s)-w(s)|\). The difference between two scores for the \textbf{same action} satisfies
\[\left|\beta\E[v(s')-w(s')\mid s,a]\right|\leq\beta\|v-w\|_\infty.\]
Since \(|\max_a x_a-\max_a y_a|\leq\max_a|x_a-y_a|\),
\[\|Tv-Tw\|_\infty\leq\beta\|v-w\|_\infty.\]
Thus \(T\) is a \term{contraction}: it reduces the maximum distance by at least the factor \(\beta<1\).
\begin{itemize}
\item There is a unique fixed point \(V^*\); every initial finite vector converges to it.
\item After \(k\) sweeps,
\[\|v_k-V^*\|_\infty\leq\beta^k\|v_0-V^*\|_\infty.\]
\item A discount factor closer to one gives a slower worst-case rate.
\end{itemize}
\end{frame}

\begin{frame}{A Stopping Rule With an Error Bound}
\small
After computing \(v_{k+1}=Tv_k\), measure the \term{Bellman residual} of \(v_k\):
\[d_k=\|Tv_k-v_k\|_\infty=\|v_{k+1}-v_k\|_\infty.\]
The contraction inequality gives
\begin{align*}
\|v_k-V^*\|_\infty
&\leq\|v_k-Tv_k\|_\infty+\|Tv_k-TV^*\|_\infty\\
&\leq d_k+\beta\|v_k-V^*\|_\infty,
\end{align*}
so \(\|v_k-V^*\|_\infty\leq d_k/(1-\beta)\).
\begin{itemize}
\item To guarantee value error at most \(\varepsilon\), stop when \(d_k\leq(1-\beta)\varepsilon\) and return \(v_k\).
\item At \(\beta=0.9\) and \(\varepsilon=0.01\), require \(d_k\leq0.001\).
\end{itemize}
\footnotesize This guarantee uses exact expectations and maximization. Numerical approximation introduces additional errors.
\end{frame}

\begin{frame}{Recovering a Policy}
For the returned value estimate \(v\), choose
\[\widehat\sigma(s)\in\argmax_{a\in\Aset(s)}
\sum_r\Pr\{R(s,a)=r\}\sum_{s'}P(s'\mid s,a,r)
\bigl[r+\beta v(s')\bigr].\]
\begin{itemize}
\item With \(v=V^*\), any maximizing action is optimal; ties are allowed.
\item With an approximate \(v\), this is a candidate policy. Evaluate it to check the return it actually achieves.
\item Greedy means maximizing \textbf{current reward plus continuation value}, not just immediate reward.
\item Recompute action scores using the returned \(v\), and break ties consistently.
\end{itemize}
\end{frame}

\begin{frame}{Example: VFI for McCall Job Search}
\small
The state is \((\mathrm U,w)\) or \((\mathrm E,w)\). Rejecting pays \(b\) and draws a new offer \(w'\) with probability \(q(w')\); accepting pays \(w\) forever.
\[v_k((\mathrm E,w))=\frac{w}{1-\beta}\quad\text{for every }k.\]
Use this known employed value and update unemployed states:
\begin{align*}
q_k((\mathrm U,w),\text{accept})&=\frac{w}{1-\beta},\\
q_k((\mathrm U,w),\text{reject})&=b+\beta\sum_{w'}q(w')v_k((\mathrm U,w')),\\
v_{k+1}((\mathrm U,w))&=\max\{q_k((\mathrm U,w),\text{accept}),\ q_k((\mathrm U,w),\text{reject})\}.
\end{align*}
\begin{itemize}
\item Start unemployed-state values at zero; keep employed-state values fixed at their exact values.
\item At convergence, accept when the acceptance score is at least the rejection score.
\end{itemize}
\footnotesize Lowercase \(q(w')\) is the offer probability; \(q_k(s,a)\) is the iteration's action score.
\end{frame}

\begin{frame}{Example: McCall Iterations}
\small
Use \(\beta=0.9\), \(b=1\), offers \((1,2,3,4,5)\), and probabilities \((0.1,0.2,0.4,0.2,0.1)\).
\begin{center}
\begin{tabular}{lrrrrr}\toprule
 & \(w=1\)&\(w=2\)&\(w=3\)&\(w=4\)&\(w=5\)\\\midrule
\(v_0\)&0.00&0.00&0.00&0.00&0.00\\
\(v_1\)&10.00&20.00&30.00&40.00&50.00\\
\(v_2\)&28.00&28.00&30.00&40.00&50.00\\
\(v_3\)&31.06&31.06&31.06&40.00&50.00\\
\(v_10\)&34.20&34.20&34.20&40.00&50.00\\
\(V^*\)&34.32&34.32&34.32&40.00&50.00\\
\bottomrule\end{tabular}
\end{center}
Columns give unemployed-state values; the last row is \(V^*\). Acceptance values are \((10,20,30,40,50)\).
\begin{itemize}
\item The first update values waiting at only \(1\); later updates incorporate the value of future offers.
\item In the limit, reject offers 1--3 and accept 4--5. For rejected offers, the common value solves
\[x=1+0.9\bigl(0.7x+0.2\times40+0.1\times50\bigr),\qquad x=\frac{12.7}{0.37}.\]
\end{itemize}
\end{frame}

\begin{frame}{Example: VFI for Continuous Saving}
\small
Stock \(s\in[0,M+B]\); save \(a\in[0,\min(s,M)]\), consume \(s-a\), and draw \(Y'\sim\mathrm{Uniform}[0,B]\) independently. Next stock is \(s'=a+Y'\). Utility is \(c^\alpha\), with \(0<\alpha<1\), \(M\geq0\), and \(B>0\).
\begin{align*}
q_k(s,a)&=(s-a)^\alpha+\frac{\beta}{B}\int_0^B v_k(a+y)\,dy,\\
v_{k+1}(s)&=\max_{a\in[0,\min(s,M)]}q_k(s,a),\\
\widehat\sigma(s)&\in\argmax_{a\in[0,\min(s,M)]}
\left\{(s-a)^\alpha+\frac{\beta}{B}\int_0^B v(a+y)\,dy\right\}.
\end{align*}
\begin{itemize}
\item Every saving choice trades current consumption against future resources.
\item The exact operator is a contraction on bounded continuous functions under the supremum norm. Rewards are bounded and feasible intervals are compact.
\end{itemize}
\end{frame}

\begin{frame}{Computing the Continuous Update}
The saving model has infinitely many states and actions. A computer needs a finite representation:
\begin{enumerate}
\item Store values at selected stock points and \term{interpolate} between them.
\item Approximate the income integral by a weighted average, using \term{numerical quadrature} or sampled incomes.
\item Maximize over the saving interval using numerical optimization or a feasible action grid.
\end{enumerate}
\par\medskip
The model remains continuous. A small iteration residual measures convergence of the numerical update; it does not by itself establish accuracy for the original model.
\par\smallskip
Refine the state representation, integration, and maximization to check that values and choices are stable.
\end{frame}

% BEGIN MANAGED READINGS: week3_vfi
\begin{frame}{References and Further Reading}
\coursereading{1}
  {Richard S. Sutton and Andrew G. Barto (2018), Reinforcement Learning: An Introduction. 2nd ed. MIT Press. [Book]}
  {https://www.andrew.cmu.edu/course/10-703/textbook/BartoSutton.pdf}
  {\begin{itemize}
\item \href{https://www.andrew.cmu.edu/course/10-703/textbook/BartoSutton.pdf\#page=80}{Section 3.5: Policies and Value Functions}
\item \href{https://www.andrew.cmu.edu/course/10-703/textbook/BartoSutton.pdf\#page=84}{Section 3.6: Optimal Policies and Optimal Value Functions}
\item \href{https://www.andrew.cmu.edu/course/10-703/textbook/BartoSutton.pdf\#page=104}{Section 4.4: Value Iteration}
\end{itemize}}
\coursereading{2}
  {Thomas J. Sargent and John Stachurski, QuantEcon. [Online lectures]}
  {https://quantecon.org/lectures/}
  {\begin{itemize}
\item \href{https://tools-techniques.quantecon.org/discrete_dp.html\#the-bellman-equation-and-the-principle-of-optimality}{Discrete State Dynamic Programming — The Bellman Equation and the Principle of Optimality}
\item \href{https://tools-techniques.quantecon.org/discrete_dp.html\#value-function-iteration}{Discrete State Dynamic Programming — Value Function Iteration}
\item \href{https://python.quantecon.org/mccall_model.html\#the-mccall-model}{Job Search I: The McCall Search Model — The McCall Model}
\item \href{https://tools-techniques.quantecon.org/discrete_dp.html\#example-a-growth-model}{Discrete State Dynamic Programming — Example: A Growth Model}
\item \href{https://dynamics.quantecon.org/ifp.html\#the-optimal-savings-problem}{The Income Fluctuation Problem I: Basic Model — The Optimal Savings Problem}
\end{itemize}}
\end{frame}
% END MANAGED READINGS: week3_vfi
\end{document}
